% Candidate observational section; not included in the active manuscript.
\section{Raw-Gradient Dynamics Along Adam Trajectories}

Adam averages and rescales gradients before applying parameter updates.
Consecutive raw-gradient reversal therefore does not directly describe reversal
of the applied updates. The raw gradients nevertheless reveal how the training
objective changes along the path selected by Adam. We examine their behavior
around pronounced, transient loss excursions.

The projected-change expression remains available:
\begin{equation}
 \eta\widehat\kappa_t
 =1-\rho_t\frac{\|g_t\|}{\|g_{t-1}\|}.
\end{equation}
Under Adam, this quantity describes a normalized projected gradient change;
it does not identify Hessian curvature or an Adam stability boundary. For
nonzero gradients, exceeding two requires negative raw cosine and gradient-norm
growth. These linked observations describe one geometric pattern, rather than
independent evidence for several mechanisms.

\subsection{Controlled Trajectories and Local Spikes}

We reuse the three controlled full-batch objectives: CIFAR-10/MLP,
CIFAR-10/ResNet-32 and WikiText-2/Transformer. Each case starts from one frozen
initialization and uses three constant Adam rates for 1,024 updates, with fresh
moments and no clipping, decay or learning-rate schedule. The overview shows
training loss, raw projected change and raw cosine across all nine arms.
The loss is unsmoothed. The other curves use 16-observation trailing means;
the projected-change panels focus their axes on
$[-0.1,2.2]$. Timing statements use raw observations, including values outside
this display range.

The spike figure shows the completed episode with the largest loss peak relative
to its baseline at each case's highest rate. All observations are unsmoothed.
For the Transformer example, projected change crosses two at state 788, when
loss is only 0.40\% above the recent minimum; loss reaches 12.0\% above that
minimum at 790 and peaks at 794. Raw cosine is negative, and gradient and
momentum norms grow before the pronounced rise. This crossing precedes the
large increase, although the first small positive loss change has already
occurred. The other cases show related local behavior, with differences in
the ordering of the norm and loss maxima.

\subsection{Recurrence and Its Limits}

A simple loss-only census checks repetition beyond these selected examples.
An episode begins when loss rises at least 10\% above the preceding 16-state
minimum and ends at its first return below that frozen baseline. Rebounds
before recovery remain one episode. The nine arms contain 129 episodes;
45 reach at least 50\% above baseline and are called large episodes.

The recurrence table retains every arm, including misses and zero-event arms.
At the highest rates, a fresh crossing of projected change above two occurs
1--8 updates before the 10\% onset, while loss is still within 5\% of its own
recent minimum, in 13/17 large MLP episodes,3/3 ResNet episodes and 7/7 Transformer
episodes. Conditional median leads are 5,2 and 5 updates. Across all episode
sizes, the corresponding counts are 17/39,5/7 and 13/13. These observations show
recurrence without suggesting that every excursion shares the same precursor.

Crossing two is also not sufficient for an excursion. The smallest-rate ResNet
and Transformer arms cross two despite having no qualifying loss episode.
The figure examples and episode counts therefore support a recurring temporal
association along these adaptive trajectories. They do not establish a causal
mechanism or a general warning rule. Repeated episodes are dependent observations
from one initialization per case, and the analysis is exploratory. The present
claim concerns their observed geometry and timing; interventions are deferred.
